% Chapter 8 worked example. Included by ch08_resources.tex.
\section{Worked Example: Imaging Prioritisation Decision Model}
\label{sec:comp-ch08-worked-example}

\subsection{Purpose and Status}
This worked example demonstrates how the methods in Chapter~8 can be
translated into a reviewable institutional decision. All numerical inputs are
synthetic unless a source is identified explicitly. They are intended to show the
logic of the analysis, not to report empirical findings or establish a universal
threshold. Local use requires replacement of the assumptions with current,
versioned evidence and approval through the relevant clinical, managerial,
economic, legal, technical, and governance processes.

A screening programme is considering AI-supported worklist prioritisation. The model changes reading order rather than replacing professional interpretation. Value depends on queue effects, false-positive and false-negative consequences, time to diagnosis, capacity, and longer-term outcomes. 

\subsection{Decision Problem and Comparators}
\textbf{Decision problem.} Should the programme introduce AI-supported prioritisation, retain the current first-in-first-out worklist, or adopt specialist manual triage?

The comparator set is deliberately broader than the existing pathway. This prevents
a new technology from receiving a lower evidentiary threshold than staffing,
workflow, shared-service, conventional digital, or no-change alternatives.

\begin{longtable}{@{}p{0.08\textwidth}p{0.84\textwidth}@{}}
\caption{Comparators for the Chapter 8 worked example}
\label{tab:comp-ch08-comparators}\\
\toprule
\textbf{Option} & \textbf{Comparator or strategy} \\
\midrule
\endfirsthead
\toprule
\textbf{Option} & \textbf{Comparator or strategy} \\
\midrule
\endhead
\bottomrule
\endlastfoot
1 & Current first-in-first-out worklist \\
2 & Specialist manual triage \\
3 & AI-supported priority queue with unchanged human interpretation \\
\end{longtable}

\subsection{Model and Decision Structure}
The analytical structure should follow the decision rather than the available
software. The team should map the causal pathway from intervention input to
information, human decision, action, outcome, resource consequence, and review.
The structure should also identify capacity constraints, uncertainty, subgroup
variation, implementation requirements, and events that would change the
recommendation.

A compact quantitative relationship used in this example is:
\begin{equation}
\mathrm{EVPI}=\mathrm{E}_{\theta}[\max_j \mathrm{NMB}_j(\theta)]-\max_j\mathrm{E}_{\theta}[\mathrm{NMB}_j(\theta)]
\label{eq:comp-ch08-key-relationship}
\end{equation}
The equation is an organizing aid rather than a substitute for disaggregated evidence,
qualitative findings, or mandatory safeguards. Its variables and interpretation must
be defined in the local protocol.

\subsection{Synthetic Parameters and Evidence Sources}
Table~\ref{tab:comp-ch08-parameters} provides the base-case inputs. A
production analysis should add parameter distributions, correlations, structural
alternatives, data dates, system versions, and evidence-confidence ratings.

\begin{longtable}{@{}p{0.32\textwidth}p{0.22\textwidth}p{0.38\textwidth}@{}}
\caption{Synthetic parameters for the Chapter 8 worked example}
\label{tab:comp-ch08-parameters}\\
\toprule
\textbf{Parameter} & \textbf{Base value} & \textbf{Source or interpretation} \\
\midrule
\endfirsthead
\toprule
\textbf{Parameter} & \textbf{Base value} & \textbf{Source or interpretation} \\
\midrule
\endhead
\bottomrule
\endlastfoot
Annual examinations & 25,000 & Synthetic worked-case input \\
Cancer prevalence & 0.8\% & Synthetic input \\
Priority sensitivity & 0.88 & Synthetic input \\
Priority specificity & 0.82 & Synthetic input \\
Mean cancer reporting time, current & 5.0 days & Queue assumption \\
Mean cancer reporting time, intervention & 1.0 day & Queue assumption \\
Annual programme cost & GBP 85,000 & Synthetic lifecycle estimate \\
QALY loss per material delay & 0.10 & Synthetic decision parameter \\
Willingness-to-pay thresholds & GBP 30,000 and GBP 50,000/QALY & Illustrative thresholds \\
\end{longtable}

\subsection{Step-by-Step Analysis}
The sequence below is intended to make the analysis reproducible and to ensure
that evidence, economics, governance, and implementation develop together.

\begin{longtable}{@{}p{0.07\textwidth}p{0.55\textwidth}p{0.30\textwidth}@{}}
\caption{Analytical sequence}
\label{tab:comp-ch08-analysis-steps}\\
\toprule
\textbf{Step} & \textbf{Required work} & \textbf{Decision record} \\
\midrule
\endfirsthead
\toprule
\textbf{Step} & \textbf{Required work} & \textbf{Decision record} \\
\midrule
\endhead
\bottomrule
\endlastfoot
1 & Specify the comparator, human reading workflow, priority threshold, queue discipline, and downstream pathway. & Evidence and responsible owner to be recorded \\
2 & Use a classification decision tree to estimate true-positive, false-negative, false-positive, and true-negative consequences. & Evidence and responsible owner to be recorded \\
3 & Use discrete-event simulation or a queueing model to estimate time-to-report distributions under constrained reader capacity. & Evidence and responsible owner to be recorded \\
4 & Link material diagnostic delay to downstream cost and QALY consequences. & Evidence and responsible owner to be recorded \\
5 & Represent subgroup prevalence, performance, access, and delay distributions explicitly. & Evidence and responsible owner to be recorded \\
6 & Conduct deterministic sensitivity analysis, probabilistic sensitivity analysis, and value-of-information analysis. & Evidence and responsible owner to be recorded \\
\end{longtable}

\subsection{Illustrative Outputs}
The outputs in Table~\ref{tab:comp-ch08-outputs} show how technical,
clinical, operational, economic, equity, and governance findings should be kept
visible rather than compressed into a single favourable or unfavourable score.

\begin{longtable}{@{}p{0.30\textwidth}p{0.24\textwidth}p{0.38\textwidth}@{}}
\caption{Illustrative model and decision outputs}
\label{tab:comp-ch08-outputs}\\
\toprule
\textbf{Output} & \textbf{Result} & \textbf{Interpretation} \\
\midrule
\endfirsthead
\toprule
\textbf{Output} & \textbf{Result} & \textbf{Interpretation} \\
\midrule
\endhead
\bottomrule
\endlastfoot
Mean cancer reporting time, comparator & 5.0 days & Synthetic \\
Mean cancer reporting time, prioritisation & 1.0 day & Synthetic \\
Materially delayed diagnoses avoided & 14.26 & Synthetic calculation \\
Incremental QALYs & 1.43 & Synthetic calculation \\
Net incremental cost & GBP 58,914 & Synthetic calculation \\
ICER & GBP 41,305/QALY & Synthetic calculation \\
Probability cost-effective at GBP 30,000/QALY & 28.8\% & Synthetic PSA \\
Probability cost-effective at GBP 50,000/QALY & 57.8\% & Synthetic PSA \\
\end{longtable}

\subsection{Sensitivity, Uncertainty, and Subgroups}
At minimum, the completed analysis should vary the parameters most likely to
change the decision, test at least one credible alternative model structure, and
report subgroup results separately where performance, access, response, burden,
or opportunity cost may differ. Deterministic analysis should identify thresholds;
probabilistic analysis should be used where credible parameter distributions and
correlations are available. Scenario analysis is preferable when structural or
future uncertainty cannot be represented defensibly by one probability model.

The record should distinguish uncertainty that can be reduced through evidence,
uncertainty that can be managed through safeguards, and uncertainty that remains
too consequential to accept. Missing evidence should not be represented as an
average or neutral value without justification.

\subsection{Interpretation and Recommendation}
Undertake a controlled local pilot rather than unrestricted scale. The decision remains sensitive to programme price, the relationship between delay and outcome, queue performance, local calibration, and subgroup effects.

The recommendation is conditional rather than permanent. It applies only to the
specified population, setting, intervention configuration, evidence cut-off, and
version. The following conditions should be incorporated into the approval,
contract, implementation, monitoring, or renewal record:

\begin{itemize}
 \item the model remains prioritisation support rather than autonomous interpretation.
 \item reader and scanner capacity assumptions are validated locally.
 \item deprioritised examinations and subgroup delay are monitored.
 \item the model, threshold, queue logic, and workflow version are archived.
 \item renewal depends on observed time-to-diagnosis, outcomes, workload, cost, and equity.
\end{itemize}

\subsection{Decision Statement}
\begin{quote}
For the specified decision problem and comparator set, the institution should
\emph{[proceed / proceed with conditions / redesign / pilot / restrict / defer /
replace / retire]}. The decision applies to \emph{[population, setting, version,
and evidence period]}, is owned by \emph{[accountable authority]}, and will be
reviewed on \emph{[date]} or earlier if \emph{[trigger events]} occur. The
evidence, assumptions, unresolved uncertainty, mandatory safeguards, monitoring
thresholds, and exit arrangements are recorded in the accompanying dossier.
\end{quote}
